\scanpage{31}
\mspage{1402}\textbf{--- Série A}\hfill\textbf{C. R. Acad. Sc. Paris, t. 278 (27 mai 1974)}\par\noindent\rule{\linewidth}{.35pt}\par

2.5. \textsc{Proposition.} --- \textit{Soient \(\mathrm{W},\mathrm{W}'\) deux surfaces à bords \(\partial\mathrm{W},\partial\mathrm{W}'\) orientées et de même genre (nombre d'anses), et \(h\) un homéomorphisme orienté de \(\partial\mathrm{W}\) sur \(\partial\mathrm{W}'\); alors \(h\) s'étend en un homéomorphisme orienté de \(\mathrm{W}\) sur \(\mathrm{W}'\).}

3. \textsc{Voisinage tubulaire standard d'un 1-complexe circularisé.} --- Dans la suite, \(\mathrm{K}\) désigne un 1-complexe compact.

3.1. On appellera \textit{1-complexe circularisé} le couple \((\mathrm{K},\omega)\) fermé de \(\mathrm{K}\) et d'une famille (circularisation) \(\omega=(\omega_x)_{x\in\mathrm{S}}\), où pour tout sommet \(x\in\mathrm{S}\), \(\omega_x\) est un ordre circulaire sur l'ensemble des arcs issus de \(x\). [Cette notion a d'abord été systématisée dans le contexte combinatoire dans (\textsuperscript{1}).]

3.2. A tout 1-complexe circularisé \((\mathrm{K},\omega)\) est associé son \textit{voisinage tubulaire standard} \(\mathrm{T}(\mathrm{K},\omega)\) qui est une surface à bord orientée obtenue par réunion et recollement de la façon suivante : à tout sommet isolé \(x\) on associe un disque orienté \(\mathrm{D}_x\); à tout arc \(a\in\mathrm{A}\) on associe le rectangle \(\bar a\times\mathrm{I}\) (où \(\mathrm{I}=[-1,+1]\)) avec l'orientation produit; (\(\bar a\) est l'arc compactifié par 2 extrémités); on identifie les points de \(\bar a\times\mathrm{I}\) et de \(\overline{\sigma(a)}\times\mathrm{I}\) selon la loi \((p,t)\sim(p,-t)\). Pour tout arc \(a\) ayant une source \(x\in\mathrm{S}\), on considère \(a'=\omega_x(a)\) et on recolle les rectangles orientés \(\bar a\times\mathrm{I}\) et \(\bar a'\times\mathrm{I}\) en identifiant les points \((x,t)\) et \((x,-t)\) pour \(t\in[0,1]\).

3.3. On a un plongement canonique \(i_\omega\) de \(\mathrm{K}\) dans \(\mathrm{T}(\mathrm{K},\omega)\), induit par le plongement canonique \((p\to(p,0))\) de \(\bar a\) dans \(\bar a\times\mathrm{I}\), un sommet isolé \(x\) étant envoyé sur le centre du disque \(\mathrm{D}_x\) associé. Le plongement \(i_\omega\) est caractérisé à équivalence près (2.4) par la propriété d'être une équivalence d'homotopie.

3.4. Soit \((\mathrm{K},\omega)\) un 1-complexe circularisé et \(\mathrm{E}\) l'ensemble des arcs et des sommets isolés \(\mathrm{K}\). On définit une permutation \(\varphi_\omega\) de \(\mathrm{E}\) de la manière suivante : soit \(a\) un arc ayant une source, soit \(x\) son but, on pose \(\varphi_\omega(a)=\omega_x(-a)\). Sur les autres éléments de \(\mathrm{E}\), \(\varphi_\omega\) est l'identité. Un \(\omega\)-\textit{circuit} de \(\mathrm{K}\) est une orbite de \(\varphi_\omega\). On note \(\mathrm{C}(\mathrm{K},\omega)\) l'ensemble des \(\omega\)-circuits de \(\mathrm{K}\). On définit une bijection canonique entre \(\pi_0(\partial\mathrm{T}(\mathrm{K},\omega))\) et \(\mathrm{C}(\mathrm{K},\omega)\).

4. \textsc{Voisinage tubulaire d'un plongement.}

4.1. \textsc{Théorème.} --- \textit{Étant donné un plongement \(f\) de \(\mathrm{K}\) dans une surface orientée \(\mathrm{V}\), il existe une unique circularisation \(\omega(f)\) de \(\mathrm{K}\) telle qu'il existe une surface à bord \(\mathrm{T}(\mathrm{K},f,\mathrm{V})\) contenue dans \(\mathrm{V}\) (appelée \textup{voisinage tubulaire de} \(f\)), l'intérieur de \(\mathrm{T}(\mathrm{K},f,\mathrm{V})\) contenant \(f(\mathrm{K})\) et telle que le plongement \(f\) de \(\mathrm{K}\) dans \(\mathrm{T}(\mathrm{K},f,\mathrm{V})\) soit équivalent au plongement canonique de \(\mathrm{K}\) dans \(\mathrm{T}(\mathrm{K},\omega(f))\). Il existe alors une bijection canonique \(\psi_f\) entre \(\pi_0(\partial\mathrm{T}(\mathrm{K},f,\mathrm{V}))\) et \(\mathrm{C}(\mathrm{K},\omega(f))\),} (3.4).

4.2. \textsc{Corollaire.} --- \textit{Soit \(f_1\) (resp. \(f_2\)) un plongement de \(\mathrm{K}\) dans \(\mathrm{V}_1\) (resp. \(\mathrm{V}_2\)); on a \(\omega(f_1)=\omega(f_2)\) si et seulement si il existe des voisinages \(\mathrm{T}_1\) et \(\mathrm{T}_2\) de \(f_1(\mathrm{K})\) et \(f_2(\mathrm{K})\) tels que les plongements induits \(g_1\) de \(\mathrm{K}\) dans \(\mathrm{T}_1\) et \(g_2\) de \(\mathrm{K}\) dans \(\mathrm{T}_2\) soient équivalents.}

5. \textsc{Théorème de classification.}

5.1. Si \(f\) est un plongement de \(\mathrm{K}\) dans la surface orientée compacte \(\mathrm{V}\), tel que \(\pi_0(f)\) soit surjectif, \(\mathrm{T}\) un voisinage tubulaire, soit \(\mathrm{V}'=\mathrm{V}\setminus\operatorname{int}\mathrm{T}\); alors \(\mathrm{V}'\) est une surface compacte orientée à bord \(\partial\mathrm{V}'=-\partial\mathrm{T}\). On construit une partition de \(\pi_0(\partial\mathrm{V}')\) en rangeant dans une même classe les éléments contenus dans une même composante connexe de \(\mathrm{V}'\). A cette classe on associe un entier positif ou nul : le genre de cette composante.
% archive sheet number: 30

